% !TEX encoding = UTF-8 Unicode
% ch02：相关工作与论文定位（中文深度解读片段）
\section{相关工作与论文定位}
\label{sec:ch02}

\begin{keybox}{先抓住本章的五件事}
\begin{itemize}
  \item EoS 的核心理论脉络（catapult $\to$ EoS 定义 $\to$ 渐进尖锐化 $\to$ 自稳定化 $\to$ 饱和定理 $\to$ 相图）几乎全部建立在\textbf{确定性 / 全批}分析之上；一旦换到随机体制（mini-batch、动量、自适应），“阈值是多少”与“是什么量在饱和”这两件事都会改变。
  \item 与本文最贴近的两条先例线是\textbf{方向曲率}（CDAT 的更新方向曲率、Kalra 的临界尖锐度 $\lambda_c$）与\textbf{投影掉顶部特征向量}（bulk 投影）；本文与它们各有一个关键的设计分叉点。
  \item 论文在学术归属上非常克制：$\lambda_c$ 的瑞利商恒等式、热启动幂迭代都明确\textbf{不}算本文贡献——本文做的是“验证 $\pm0.5\%$，再重新定位为监测器”。
  \item SAM 家族（含 ASAM、SAF）属于\textbf{全局动作}；本文的负结果只针对秩-1 局部动作，不外推到它们。
  \item “预测器 / 执行器”二分是本文独创的组织轴：它把上面所有文献重新排了队，也解释了为什么“便宜地测”成立而“便宜地操控”失败。
\end{itemize}
\end{keybox}

\noindent 本章对应论文第 2 节（Related work and
positioning）。如果你已经读过教程，那么你已经知道 EoS 现象本身、$2/\eta$
阈值与自稳定化的直觉图像；本章的任务不是再讲一遍这些现象，而是回答一个更上层的问题：\textbf{这篇论文站在一条已经相当拥挤的研究线的哪个位置？它从每条先例线里各拿了什么、又在哪里分道扬镳？}理解这一点非常重要，因为论文的核心命题——“代理量是预测器而非执行器”——只有在看清前人各自追求过什么目标之后，才能显出它是一个新的、干净的二分，而不是若干已知结论的重新包装。

\subsection{EoS 核心脉络：从 catapult 到相图}
\label{sec:ch02-core}

按时间顺序把主线串一遍（教程里讲过现象，这里补的是“谁贡献了哪块理论砖”）：

\textbf{catapult 阶段。} Lewkowycz 等 (2020) 确立了大学习率阶段的 catapult（弹射——大学习率训练早期，参数沿海森最尖方向被“弹射”离开、直到该方向曲率降到稳定阈值之下，此后训练的稳定性由海森谱决定，故名弹射）
机制：在这个阶段里，\textbf{海森谱本身决定稳定性极限}（谱 = 矩阵全部特征值的集合，“海森谱”即每个参数方向上损失曲率的完整清单）。这是“用谱语言谈论训练稳定性”这一思路的起点——后来的
EoS、$2/\eta$ 阈值乃至本文的“自稳定帽”，用的都是这套词汇。

\textbf{EoS 的定义。} Cohen 等 (2021) 把现象钉死：$\lmax$ 爬升到 $\approx 2/\eta$
并在那里振荡，损失非单调下降。教程的绝大部分图像来自这篇。

\textbf{渐进尖锐化与范数代理。} Li 等 (2022) 引入渐进尖锐化（progressive
sharpening）概念，并提出了\textbf{输出层权重范数}作为尖锐度代理——这正是教程第 9
节“权重范数代理”的出处。本文对它的细化有二：\emph{last-hidden
层}是更强的跟踪者（并非输出层最优），且该代理在 BatchNorm
下\emph{失效}（相关由正转负）。这两个细化在解读第 \ref{sec:ch04-tier0}
节会展开成完整数字。

\textbf{自稳定化。} Damian 等 (2023) 证明了自稳定化机制，核心是三次恒等式 $\nabla S
= \nabla^{3}L(u,u)$（左边是尖锐度 $S=\lmax$ 的梯度；右边是损失三阶导数张量——一个“三维数字方块”——把其中两个槽位填上顶部特征向量 $u$（本文记 $\vtope$/$u_1$）做内积后缩并成的一个向量；恒等式即“尖锐度的梯度 = 曲率沿顶部方向的变化率”）：在 EoS
振荡中，沿顶部特征方向的曲率变化率被三阶项锁定。这是全文最重要的单块理论砖——本文的“自稳定帽”理论（$2/\eta$
是被周期-2
振荡\emph{顶住}的帽子，局部阻尼只能留下帽子或摘掉帽子，不能把帽子压低）正是对这套理论的延伸。

\textbf{饱和定理。} Agarwala 等 (2022) 在二阶回归模型（输出线性、
损失是输出的二次函数的回归设置——曲率可解析算出的“玩具损失面”）上给出 $2/\eta$
饱和定理，并注意到 $\lambda_2$
保持在边界\emph{之下}。这个“谱缺口”（spectral gap，最大与第二大特征值之差
$\lambda_1-\lambda_2$）观察与本文有两层关系：它佐证了 EoS
几何默认是“一维尖锐”的；同时，当本文在构造的多尖锐景观（$\lambda_1\!\approx\!\lambda_2\!=\!2/\eta$）上阻尼
$u_1$ 时，观察到 $\lambda_2$
\emph{纹丝不动、没有谱流}——这恰恰反驳了有人可能从他们的谱缺口里读出的“谱迁移”图像。

\textbf{相图。} Kalra 等 (2023) 画出早期训练动力学的四区域相图（phase diagram，“相”= 参数空间中训练行为模式相同的区域，类比水的固/液/气三态；相图即标出“哪片学习率×深度×宽度区域对应哪种动力学”的地图）。其中的 sharpness-reduction 相，正是本文“精确阻尼摘掉帽子后 $\lmax$ 离开
$2/\eta$”所对应的现象学——只不过本文把它放进了自稳定帽的机制框架里，并区分了“帽被摘掉后的上逃逸”与“帽未被摘掉时的钉在原处”两种结局。

\textbf{简化模型上的早期理论。} Chen 等 (2023) 用两步（周期-2）映射分析 EoS；Ahn 等
(2023) 证明两层阈值神经元网络上的 GD 可证地展现 EoS；Song 等 (2023)
给出分岔论解释（bifurcation theory，研究参数——如学习率——连续微小变化如何令系统行为在临界值处突然换模式的理论）——EoS 轨迹对齐到周期-2 中心流形上（center manifold，系统里既不快速收缩也不快速发散的“中性”方向构成的低维集合，轨迹常被吸在上面慢慢演化；这句话的意思是训练被困在这个低维集合附近做周期-2 来回振荡）。这三篇共同把“周期-2 振荡是 EoS
的骨架”这件事夯实，而周期-2 振荡正是本文帽机制的持帽者。

\textbf{关键限定：以上全部是确定性 / 全批分析。}
这句话是论文原节的着重号，也是读者最容易带着教程直觉踩空的地方。教程里的 $2/\eta$
是全批 GD 的定律；而随机体制在质上不同：Andreyev 等 (2024) 证明 mini-batch SGD 满足
$2/\eta$ 的是\textbf{批尖锐度}（batch sharpness）而非全批 $\lmax$，并且这会把全批
$\lmax$ 压到\emph{严格低于} $2/\eta$；Andreyev 等 (2026)
把这一分析扩展到动量；对自适应方法，Cohen 等 (2022)
表明饱和的是\textbf{预条件化的}（preconditioned）顶部特征值（指在优化器自己的坐标缩放下测得的曲率——Adam 的逐坐标缩放就是一种预条件，相当于在拉伸过的参数空间里做梯度下降），饱和在优化器自己的稳定阈值上（例如 Adam
在 $\beta_1{=}0.9$ 时是 $38/\eta$；数值上恰与后文 ResNet-18 的重球边界 38 相同，但来自不同公式，纯属巧合）。本文的秩-1 plant
可以视为这个“预条件尖锐度均衡”的秩-1
类比。同时要强调：本文的自稳定帽机制依托的是确定性的 Damian 等理论，因此\emph{在全批
GD 体制下是精确的}；动量与 mini-batch
的结果只做定性处理，并对照重球动量边界（见解读第 \ref{sec:ch03-threshold}
节）。这就是论文把“机制随体制而变、结论不变”写成显式主题的原因。

\subsection{方向曲率：最近的先例与一次关键分叉}
\label{sec:ch02-directional}

这一小节涉及与本文预测端最贴近的两篇工作，值得放慢速度。

\textbf{CDAT。}（CDAT 是该方法的名字，Roulet 等 2024 论文提出的
曲率感知学习率调节器。）Roulet 等 (2024) 的 CDAT 使用沿\textbf{更新方向} $u$ 的方向曲率（directional curvature，损失面沿给定方向的弯曲程度）$d_t
= |u^{\top}\nabla^{2} f\, u|$ 作为曲率感知的步长分母，并且\emph{拥抱} $2/\eta$
均衡（``均衡''即 EoS 轨道长期驻留的状态，第九章给出其形式化）——它不试图逃离 EoS，而是顺着 EoS
调步长。它的目标是步长\textbf{调谐}（tuning），不是 $\lmax$ \textbf{监测}，更不是
EoS 逃逸（这两个目标它根本不追求，所以论文也不在这两个维度上评价它）。本文与 CDAT
的分叉是正交的，也是本文贡献列表的第二条：\textbf{调谐}想要的是（预条件）更新方向上的曲率；\textbf{监测}
$\lmax$ 想要的则是最好的探针方向——论文证明那是原始梯度
$\mathbf{g}$，而不是动量平滑后的更新（解读第 \ref{sec:ch04-probe} 节给出不重叠 CI
的证据）。一句话：同是“方向曲率”，问的问题不同，答案的方向就不同。

\textbf{临界尖锐度。} Kalra 等 (2026)（注：本文引用的 Kalra 等
2023/2024/2026 为同一第一作者 Kalra 的不同论文）定义临界尖锐度（critical sharpness，沿更新方向线搜定义的尖锐度）$\lambda_c =
2/\eta_c$（$\eta_c$ 即线搜找到的``恰好还能保持稳定''的临界步长：步长超过它、沿该方向的更新就开始发散，$\lambda_c=2/\eta_c$ 即该方向对应的稳定阈值曲率），通过沿更新方向的线搜（line search，沿一个固定方向试几个不同步长、各做一次前向、挑损失最小的——经典优化里调步长的做法，现代深度学习训练几乎不用；5--6 次前向传播，规模推到 7B 模型），并在其结果 2.1
中陈述了恒等式 $\lambda_c =
\mathbf{g}^{\top}H\mathbf{g}/\mathbf{g}^{\top}\mathbf{g}$ 以及 $\lambda_c \le \lmax$
的界。论文在这里的表态值得逐字体会：\textbf{本文不把这个恒等式据为新颖}——它做的事是在二次型上把恒等式验证到
$\pm0.5\%$，然后把这个量从“步长调谐工具”\textbf{重新定位}为“单次 HVP 的 $\lmax$
监测器”，并刻画其适用域（梯度与顶部特征向量对齐的 EoS
体制）。这是一种很干净的增量方式：数学是别人的，验证与再定位是自己的，而“再定位”恰恰打开了本文后半的控制端问题——一个如此好的预测器能否反过来当执行器？

\subsection{投影掉顶部特征向量：bulk 投影的再解释}
\label{sec:ch02-bulk}

这是与本文\textbf{控制端}（plant）最近的先例。Song 等 (2025) 表明：把 SGD
更新投影到谱的 bulk（spectral bulk，谱主体——去掉 top-$k$
特征向量分量后的其余部分）上，训练在损失与精度上\emph{完全不变}，但尖锐度\emph{无约束增长}。这个现象正是本文“体制
(ii)”——精确阻尼解除自稳定帽——的定性症状。

本文与它的区别有两处，且都值得细品：

\begin{enumerate}
  \item \textbf{解释的层次不同。} Song 等 (2025)
 把结果读作“梯度与顶部特征向量的对齐是‘伪’对齐”，本质上是一个泛化论证；本文给出的则是\textbf{稳定性机制}：他们的投影恰好是本文秩-1
 阻尼在 $\gamma\!\to\!1$（$\gamma$ 此处为阻尼强度 $=1-1/d$，见第九章；与 BatchNorm 的仿射参数 $\gamma$ 无关——$\gamma\to1$ 即 $d\to\infty$、完全投影）下的极限，而本文理论的推论给出阻尼后均衡尖锐度 $S_{*} =
2d/\eta$——他们观察到的“无约束增长”，正是这个均衡随 $d$
 走向无穷的发散。同一现象，从“对齐是假的”变成“帽子被彻底摘了”。
  \item \textbf{设计意图不同。} 本文的 plant
刻意使用\textbf{部分阻尼}（$\gamma<1$）而不是完全投影，目的正是\emph{诊断}帽子（阻尼轻一点，看帽子还在不在），而非\emph{移除}帽子。完全投影摘帽之后什么都能长，反而看不出帽的存在。
\end{enumerate}

\subsection{精确测量的顶端}
\label{sec:ch02-exact}

Ramasinghe 等 (2026) 展示了热启动的幂迭代——以上一步的特征向量为种子（经 Adam
预条件子传输：即先把上一步的特征向量按优化器的坐标缩放变换一下、再当迭代种子）——能在 3B 规模以每步 $\lesssim 5$ 次 HVP 跟踪精确 $\lmax$，并在
Lipschitz-Hessian 加特征间隙假设下（Lipschitz 海森 = 假设曲率随参数的变化速率有上界、“曲率地图是平滑的”；特征间隙 = $\lambda_1-\lambda_2$，前两名特征值之差）给出 Davis--Kahan
漂移界（Davis–Kahan 是经典矩阵扰动定理：特征间隙越大、扰动越小，特征向量被带偏的角度就越小——大意是“前两名特征值差得够远时，顶部方向对噪声稳健”；漂移界即逐步训练中该偏差随时间累计的上界）。本文直接把它采纳为代理阶梯的\textbf{精确顶层}与 ground-truth 测量（实测中位
4--11
HVPs/步），并同样明确声明：热启动\textbf{不是}本文贡献。换言之，本文的“阶梯”顶层是借来的尺子——但正因为尺子经过校准，下方所有廉价代理才能得到严肃的评价（解读第
\ref{sec:ch03-curvature} 节讲这把尺子本身如何被校准）。

\subsection{归一化、平坦度与 SAM 家族}
\label{sec:ch02-flatness}

平坦度--泛化联系是这条线的经典动机：Keskar 等 (2017)
观察到大批量训练找到更尖锐、泛化更差的极小值；Dinh 等 (2017)
给出著名告诫——尖锐的极小值也可能泛化得很好。归一化层面，Wei 等 (2019) 与 Lyu 等
(2022) 研究了归一化层对尖锐度的影响（分别针对 Fisher
矩阵——与海森同族的“曲率/敏感度”矩阵，自然梯度与 K-FAC 用的就是它——与球面尖锐度——在当前参数周围画一个固定半径小球、看球面上损失变化有多剧烈的一类平坦度定义，而非只看单个方向的二阶曲率）。本文在此的增量是\textbf{时间分辨}的 $\|\gamma\|$ 与 $\lmax$
的相关及其随归一化族（plain / BN /
LN）的依赖关系——把“归一化影响平坦度”从静态命题变成训练全程可跟踪的监测地图（解读第
\ref{sec:ch04-tier0} 节）。

控制端方面，Foret 等 (2021) 的 SAM 是本文的逃逸基准。尤其值得留意 Long 等
(2024)：他们推导出 SAM 专属的 EoS 截断（cutoff，此处指 SAM 情形下被压低的稳定阈值上限）——依赖 $\|\nabla L\|$，与纯 GD 的 $2/\eta$
不同——且经验上 SAM 就跟踪这条截断。这就回答了“SAM 为什么坐在 $2/\eta$
之下”，与本文“秩-1
局部动作为什么压不下去”的自稳定帽解释\emph{互补}：一个说明全局动作为何能下去，一个说明局部动作为何下不去，两条合起来才是完整的二分。

最后是范围限定：本文测试的是规范的全局动作基线（vanilla SAM）；更广的
sharpness-aware 家族——尺度不变的扰动（scale-invariant，扰动效果不受参数整体缩放影响的：SAM 的扰动效果会随参数尺度改变，Kwon 等 2021 的 ASAM 改用按参数幅度归一的扰动来消除这一依赖）、输出层面的轨迹损失（Du
等 2022 的 SAF）——\textbf{共享全局动作结构}，因此本文的秩-1
负结果\textbf{不能}外推到它们（头对头比较留作未来工作）。读负结果论文时，这个边界划在哪里至关重要。

\subsection{论文的切位：二分作为组织轴}
\label{sec:ch02-position}

把上面的线摆在一起，就能看清本文的“切位”。前人的工作可以按“想拿尖锐度做什么”分成几类：有人解释现象（Cohen、Damian、Agarwala、Kalra
的相图、Chen / Ahn / Song 的简化模型理论），有人测量它（Kalra 的
$\lambda_c$、Ramasinghe 的热启动幂迭代），有人顺着它调步长（CDAT），有人投影它（Song
的 bulk 投影），有人绕开它找更平的极小值（SAM
家族）。\textbf{没有一篇把“预测”与“控制”当作两个必须分开回答的问题。}本文的独创之处正在于此：它借用控制论的观测器--控制器分离原理（observer--controller separation principle）作类比。这条原理是线性控制论的精确定理：\emph{观测器}（observer，从带噪输出估计系统内部状态的模块——只估计、不干预）与\emph{控制器}（controller，决定施加多大干预的决策逻辑）可以分开设计、拼起来仍然稳定。它之所以值得命名为“原理”，是因为这件事并不平凡——直觉上两者互相纠缠（控制依赖估计，估计又被控制扰动），而线性系统里这种纠缠恰好可以解除。论文借它类比“测尖锐度（观测）”与“控尖锐度（控制）”是两个可分开回答的问题，并诚实声明这只是类比、不是定理——分离原理对线性系统是精确的，而神经网络优化是非线性、非凸的，在
ResNet
实验里还是随机且带动量的，连阈值本身都会随体制移动——线性世界里``估计不被控制干扰''的保证在此不再自动成立（解读第七章的“自致传感失败”——执行器亲手毁掉自己的传感器——正是这种纠缠在非线性世界现身的实例）。论文因此把“便宜地测尖锐度”与“便宜地操控尖锐度”拆成两个独立的问题，对每个问题给出带多种子置信区间的答案：测，行；操控，不行。

\begin{keybox}{这一节与后续解读的接口}
\begin{itemize}
  \item 阈值随体制变化（全批 $2/\eta$ / 重球 $2(1{+}\mu)/\eta$ / 批尖锐度）$\Rightarrow$ 解读第 \ref{sec:ch03-threshold} 节。
  \item $\lambda_c$ 恒等式 + $\pm0.5\%$ 验证 + 重新定位为监测器 $\Rightarrow$ 解读第 \ref{sec:ch04-rayleigh} 节。
  \item bulk 投影 = $\gamma\to1$ 极限、部分阻尼诊断帽 $\Rightarrow$ 论文控制端（解读后文）。
  \item SAM 截断与帽机制的互补 $\Rightarrow$ 论文讨论部分。
\end{itemize}
\end{keybox}

\subsection*{本章要点}
\begin{itemize}
  \item EoS 主线文献（catapult、EoS
定义、渐进尖锐化、自稳定化、饱和定理、相图、简化模型理论）\textbf{全部}是确定性 /
全批分析；随机体制下饱和的量与阈值都不同（批尖锐度、预条件特征值、$38/\eta$），这决定了本文必须按体制分别陈述机制。
  \item 与预测端最近的两篇先例各有一个分叉：CDAT
用\emph{更新}方向调步长，本文证明\emph{原始梯度}才是 $\lmax$
监测的更好探针；$\lambda_c$ 恒等式属于 Kalra 等 (2026)，本文只做 $\pm0.5\%$
验证并把它重新定位为单 HVP 监测器。
  \item 与控制端最近的先例是 bulk 投影：它是本文秩-1 阻尼的 $\gamma\to1$
极限，其“尖锐度无约束增长”被本文推论 $S_{*}=2d/\eta$ 解释为该均衡的发散；本文 plant
用 $\gamma<1$ 的部分阻尼来\emph{诊断}帽子。
  \item 精确测量顶端（热启动幂迭代）借自 Ramasinghe 等 (2026)，作为 ground truth
而非贡献。
  \item “预测器 / 执行器”二分是本文独创的组织轴；SAM 家族的全局动作不在秩-1
负结果的外推范围内。
\end{itemize}
